\subsection{Objetivos}
Desarrollar un asistente conversacional que permita a los usuarios realizar consultas en español y recibir recomendaciones de películas mediante la implementacion y comparacion de dos metodos de recuperacion de informacion, uno basado en TF-IDF y otro con embeddings multilingues, ambos con similitud del coseno.

Los objetivos especi­ficos son:
\begin{enumerate}
\item Preparar un conjunto de datos de peli­culas que incluya informaciión textual relevante, como generos y sinopsis, para utilizarlo en la recuperacio­n de contenidos a partir de consultas en español.
\item Implementar un metodo de recuperacion basado en TF-IDF y similitud del coseno que permita comparar las consultas de los usuarios con la informacion textual de las peli­culas y recuperar los resultados mas relacionados.
\item Implementar un metodo de recuperacion semantica mediante embeddings multilingues y similitud del coseno, con el fin de identificar peli­culas relacionadas con la intencion de las consultas, incluso cuando existan diferencias entre los terminos utilizados por el usuario y los textos del catalogo.
\item Evaluar y comparar los dos metodos de recuperacion mediante las metricas Top-1, Recall@5, rango reciproco medio (MRR) y tiempo de respuesta, utilizando un conjunto de consultas de prueba y criterios de evaluacion definidos previamente.
\end{enumerate}
